\chapter*{Reading assessments in context}
\addcontentsline{toc}{chapter}{Introduction: Reading assessments in context}

An audit takes a view of the system at a particular point in time. It may find a defect, identify a risk, measure a weakness in the evidence or propose a different direction. Those are distinct kinds of finding. The value of the assessment depends on the question it asked, the source revision it examined and the observations it retained.

This part gathers quality and readiness assessments, performance notes, archived explorations and discarded directions. Their presence preserves a line of inquiry. It does not assign each finding the same confidence or imply that every recommendation was adopted.

\section*{Follow the evidence boundary}

For a defect claim, look for the trigger, the observed result and a check that can distinguish failure from success. For a performance claim, look for the workload, environment, measurement method and limits. For a recommendation, identify the need it serves and the assumptions it requires.

A static reading can identify a plausible risk without reproducing a runtime failure. A benchmark can establish a result for its recorded workload without establishing a general performance guarantee. A later implementation can resolve a finding while leaving the original audit useful as the explanation of how it was discovered.

\section*{Use the archive as a record of choices}

Explorations and graveyard entries can explain why a tempting approach was abandoned or which missing constraint prevented adoption. Read them alongside later design records and release evidence before treating an older conclusion as current.

When applying an assessment to today's code, first establish whether the relevant path and behaviour still exist. Then reproduce or re-evaluate the finding at the revision being considered. The archive helps frame that investigation by preserving the original observation and its context.
